\section{Sensitivity of the headline result}\label{app:sensitivity}

This appendix documents the sensitivity of the empirical-comparison ranking in Section~\ref{sec:empirical_comparison} to the regulatory-uncertainty channels. The headline result is the regime classification (dominated, wins today, absorbed by Nexus) reported in Sections~\ref{sec:empirical_comparison} and~\ref{sec:discussion}. The counterparty-parameter sensitivity, the calibration residuals, descriptive evidence on stablecoin pricing around recent regulatory enactments, and a two-class robustness extension are reported in a separate internet appendix.

\subsection{Regulatory uncertainty band}\label{sub:regulatory_uncertainty_band}

The headline decomposition reported in Section~\ref{sec:identification} prices the structural counterparty and depeg components at the calibrated values of $(h, \xi)$ and the empirical stress-regime frequency $p_s$. Two recent contributions document distinct channels through which regulatory uncertainty transmits to those components transitorily, without altering the structural calibration. We use their estimates to construct an adverse band around the central expected loss and to report the implied movement under episode conditions comparable to the events both papers identify.

\paragraph{Issuer counterparty channel.} \citet{cerutti_firat_hengge_sagawa_2026} identify exogenous stablecoin demand shocks through narrative selection of regulatory and idiosyncratic events from 2019 to mid-2025 and report high-frequency equity responses for ten firms with potential stablecoin exposure. A one percent shock to the combined USDC and USDT market capitalisation raises Coinbase equity by 1.44 percent, PayPal by 0.52 percent, Square by 0.59 percent, and Adyen by 0.58 percent, with all four coefficients significant at conventional levels; equity responses for large banks, community banks, retailers, and the two card networks are statistically insignificant. The asymmetry is informative for our framework: markets price the upside of stablecoin adoption for firms with stablecoin infrastructure and do not price disintermediation risk for incumbent deposit-taking institutions, consistent with the structural separation between $\Pi_{\text{risk}}$, borne by the stablecoin holder, and the observable basis, which absorbs jurisdiction-specific compliance cost. For our purposes the relevant magnitude is the inverse mapping: an adverse regulatory shock that contracts combined market capitalisation by three to five percent maps to a five to seven percent equity correction at a stablecoin-linked issuer such as Circle or Coinbase. Under a Merton-style equity-to-default translation, an equity correction of that order raises the one-year hazard $h$ for the issuer by roughly 50 to 100 basis points. The structural counterparty component $\ell_C = h \cdot \xi \cdot 10^4$ then moves from 20 basis points at the calibrated $h = 0.01$ to a range of 30 to 40 basis points at $h \in [0.015, 0.020]$ during the adverse episode. For the issuer with the larger hazard base we apply $h \in [0.025, 0.035]$, a rise of 50 to 150 basis points reflecting the wider adverse cone that attestation-only disclosure implies for the equity translation, which moves $\ell_C$ from 60 to between 75 and 105 basis points. Persistence is governed by the half-life of the equity correction and is in practice weeks rather than days. The same mapping reads in reverse as a check on the central level of $h$ used in Section~\ref{sec:risk_factors}: at the elasticity that converts a five to seven percent equity correction into a 50-to-100 basis point hazard rise, a central hazard materially above the calibrated one percent would require a persistently depressed equity valuation of the listed issuer in calm conditions, so the equity market provides an observable that disciplines the level and not only the adverse band.

\paragraph{Stress-regime probability channel.} \citet{alasadi_bewaji_gugnani_gupta_heijmans_2025} estimate a GARCH, SVAR, and TVP-VAR framework on USDT, USDC, DAI, and TUSD volatility with the \citet{baker_bloom_davis_2016} economic policy uncertainty index entered as a structural shock. A one standard deviation shock to economic policy uncertainty raises USDC and TUSD return volatility by approximately half a percent at impact, with USDT and DAI essentially unresponsive. The time-varying connectedness measures locate economic policy uncertainty as a net volatility transmitter to stablecoins, with net pairwise spillovers to USDT peaking between 20 and 40 during the COVID-19 dislocation, the mid-2022 monetary tightening, and the March 2023 banking sector episode. Frequency-domain decomposition shows that long-term integration becomes the primary component after mid-2021, indicating that uncertainty transmission is persistent at the regime level even though single-shock responses dissipate in one to five days. Mapped to the global-game object, this transmission does not move the structural microstructure block; it shifts the unconditional probability $p_s$ of the stress regime upward for the duration of the episode of heightened uncertainty. At the endogenous $p_s = 0.0052$, the probability-weighted stress component $\ell_{D,s}$ contributes about 1.3 basis points to USDC's holding loss. Under the elevated regime $p_s \in [0.02, 0.03]$ that the connectedness peaks suggest, the same component rises to between 4.8 and 7.1 basis points for USDC, and to between 0.5 and 0.8 for USDT, whose stress depeg is far milder.

\paragraph{The combined band.} Table~\ref{tab:reg_uncertainty_band} reports the implied total expected loss under the two channels operating jointly. The baseline depeg component $\ell_{D,b}$, read from the observable off-ramp basis, is unchanged across columns. The adverse band reflects the two transitory channels: $h$ in the counterparty component and $p_s$ in the stress-regime component.

\begin{table}[H]
\centering
\caption{Holding-leg expected loss under the central parametrisation and under the adverse regulatory-uncertainty regime that combines the counterparty and stress channels. Component rows report each component's range over the adverse region, while the total row pairs the two adverse end points $(h_{\text{lo}}, p_{s,\text{lo}})$ and $(h_{\text{hi}}, p_{s,\text{hi}})$, so the rows need not sum within a column. The conversion basis is a circuit fee and is not part of the holding loss.}\label{tab:reg_uncertainty_band}
\begin{tabular}{lrrrr}
\toprule
Component (bps)             & USDC central & USDC adverse & USDT central & USDT adverse \\
\midrule
$\ell_C$                    & 20.0  & 30.0 to 40.0  & 60.0  & 75.0 to 105.0 \\
$\ell_{D,s}$ stress         & \phantom{0}1.3 & \phantom{0}4.8 to 7.1 & \phantom{0}0.0 & 0.5 to 0.8 \\
\midrule
$\Pi_{\text{risk}}$         & 21.3  & 34.8 to 47.1  & 60.0  & 75.5 to 105.8 \\
\bottomrule
\end{tabular}
\end{table}

The USDC holding loss rises by roughly 14 to 26 basis points above the central value under the adverse regime; the USDT loss by 15 to 46. The USDT band is wider in absolute terms because the counterparty component is three times the USDC one, so the same proportional shock to $h$ moves the loss by a correspondingly larger amount. None of the use-case rankings of Section~\ref{sec:empirical_comparison} flips at the upper bound of either band: the EU-Brazil margin of USDC against SWIFT is 382 basis points before the band and 356 at the upper bound, and the USDT margin is 344 before the band and 298 at the upper bound. The band characterises an episode-conditional widening of the holding loss that the regime classification in Section~\ref{sec:discussion} absorbs without reordering.
